\documentclass{beamer}

\usepackage{tikz}
\usetikzlibrary{positioning,matrix, arrows.meta}
\usepackage{color}
\newcommand{\vet}[1]{\foreach \num in {#1}{\el{\num}}}
\newcommand{\vetempty}[1]{\foreach \num in {#1}{\elempty{\num}}}
\newcommand{\el}[1]{\tikz{\node[font=\huge, minimum size=1cm, draw]{#1}}}
\newcommand{\elempty}[1]{\tikz{\node[font=\footnotesize, minimum size=1cm, draw=white, fill=white]{#1}}}
\newcommand{\bl}{\color{blue}}
\newcommand{\gr}{\color{gray}}
\newcommand{\re}{\color{red}}
\newcommand{\tcb}[1]{\textcolor{blue}{#1}}
\newcommand{\tcg}[1]{\textcolor{green}{#1}}
\newcommand{\tcr}[1]{\textcolor{red}{#1}}
\newcommand{\ar}{\tikz{\draw [->] (0,0) -- (0,0.4)}}
\newcommand{\nd}{\color{white!0}.}

\usepackage{beamerthemeSingapore}
\usepackage{graphicx}
%\usepackage{here}
\usepackage{url}
\usepackage[brazil]{babel}
\usepackage[T1]{fontenc}
\usepackage[utf8]{inputenc}
\usepackage[tight,raggedright]{subfigure}
\usepackage[alf]{abntex2cite}
\usepackage{listings}
\usepackage{minted}
\usepackage{amsmath}
\usepackage{pdfpages}
% Please add the following required packages to your document preamble:
\usepackage{booktabs}
\usepackage{ragged2e}
\usepackage{etoolbox}
\usepackage{mathtools}
\DeclarePairedDelimiter{\ceil}{\lceil}{\rceil}

\apptocmd{\frame}{}{\justifying}{} % Allow optional arguments after frame.

% TODO TODO TODO TODO TODO TODO TODO
\title{04 - Busca}
\author{Prof. Alex Kutzke}
\date{Maio de 2021}

\begin{document}

	% ==========================================
	\begin{frame}
		\begin{center}
			% TODO TODO TODO TODO TODO TODO
			\large{Setor de Educação Profissional e Tecnológica}\\
			Tecnologia em Análise e Desenvolvimento de Sistemas\\
			\vspace{1cm}			
			\small{DS141 - Estruturas de Dados II}
		\end{center}
		\titlepage
		\begin{center}
			Slides adaptados do material produzido\\pelo Prof. Rodrigo Minetto
		\end{center}
	\end{frame}
	% ==========================================

	% ==========================================
	\begin{frame}[c]{Roteiro}
		\tableofcontents

	\end{frame}
	% ==========================================

	\section{Introdução}

	\begin{frame}[c]{Busca}
		\textbf{Motivação}: a operação de \tcr{busca} por uma \tcb{palavra-chave}
		ou \tcb{valor} é a base de muitas aplicações da computação, seja em um 
		motor de busca na internet ou em um sistema comercial.
	\end{frame}

	\begin{frame}[c]{Busca}
		\begin{description}
			\justifying
			\item[O problema]: uma \tcb{coleção} de elementos, onde cada elemento possui
				um identificador e uma \tcb{chave} para busca;
			\item[Saída]: o \tcb{índice} na coleção que possui a mesma \tcb{chave} procurada
				\tcr{ou} identificar que tal chave \tcb{não existe} na coleção.
		\end{description}
	\end{frame}

	\begin{frame}[c]{Busca}
		Em nossos exemplos usaremos um vetor de inteiros como coleção. No entanto,
		podemos utilizar qualquer coleção cujos elementos possam ser comparados.
	\end{frame}

	\section{Busca sequencial}

	\begin{frame}[c]{Busca}
		A forma mais simples de fazer uma busca em um vetor consistem em \tcb{percorrer o vetor},
		\tcr{elemento a elemento}, para verificar se o elemento de interesse é igual
		a um dos elementos do vetor.
	\end{frame}

	% ==========================================
	\begin{frame}[c]{Exemplo Busca sequencial}
		\begin{itemize}
			\item \tcb{Elemento desejado:} 66.
		\end{itemize}
		\begin{center}
			\vetempty{0,1,2,3,4,5,6,7}\\
			\vet{22,99,66,33,11,00,77,55}\\
		\end{center}
		\begin{itemize}
			\item \tcb{Comparação}
		\end{itemize}
	\end{frame}
	% ==========================================
	
	% ==========================================
	\begin{frame}[c]{Exemplo Busca sequencial}
		\begin{itemize}
			\item \tcb{Elemento desejado:} 66.
		\end{itemize}
		\begin{center}
			\vetempty{0,1,2,3,4,5,6,7}\\
			\vet{\bf \bl 22,99,66,33,11,00,77,55}\\
		\end{center}
		\begin{itemize}
			\item \tcb{Comparação}
		\end{itemize}
	\end{frame}
	% ==========================================

	% ==========================================
	\begin{frame}[c]{Exemplo Busca sequencial}
		\begin{itemize}
			\item \tcb{Elemento desejado:} 66.
		\end{itemize}
		\begin{center}
			\vetempty{0,1,2,3,4,5,6,7}\\
			\vet{22,\bf \bl 99,66,33,11,00,77,55}\\
		\end{center}
		\begin{itemize}
			\item \tcb{Comparação}
		\end{itemize}
	\end{frame}
	% ==========================================

	% ==========================================
	\begin{frame}[c]{Exemplo Busca sequencial}
		\begin{itemize}
			\item \tcb{Elemento desejado:} 66.
		\end{itemize}
		\begin{center}
			\vetempty{0,1,2,3,4,5,6,7}\\
			\vet{22,99,\bf \bl 66,33,11,00,77,55}\\
		\end{center}
		\begin{itemize}
			\item \tcb{Comparação}
		\end{itemize}
	\end{frame}
	% ==========================================
	
	% ==========================================
	\begin{frame}[c]{Exemplo Busca sequencial}
		\begin{itemize}
			\item \tcb{Elemento desejado:} 66.
		\end{itemize}
		\begin{center}
			\vetempty{0,1,2,3,4,5,6,7}\\
			\vet{22,99,\bf \re 66,33,11,00,77,55}\\
		\end{center}
		\begin{itemize}
			\item \tcb{Comparação};
			\item \tcr{Elemento encontrado na posição 3}.
		\end{itemize}
	\end{frame}
	% ==========================================

	% ==========================================
	\begin{frame}[c]{Exemplo Busca sequencial}
		\begin{itemize}
			\item \tcb{Elemento desejado:} 44.
		\end{itemize}
		\begin{center}
			\vetempty{0,1,2,3,4,5,6,7}\\
			\vet{22,99,66,33,11,00,77,55}\\
		\end{center}
		\begin{itemize}
			\item \tcb{Comparação};
		\end{itemize}
	\end{frame}
	% ==========================================


	% ==========================================
	\begin{frame}[c]{Exemplo Busca sequencial}
		\begin{itemize}
			\item \tcb{Elemento desejado:} 44.
		\end{itemize}
		\begin{center}
			\vetempty{0,1,2,3,4,5,6,7}\\
			\vet{\bf \bl 22,99,66,33,11,00,77,55}\\
		\end{center}
		\begin{itemize}
			\item \tcb{Comparação};
		\end{itemize}
	\end{frame}
	% ==========================================

	% ==========================================
	\begin{frame}[c]{Exemplo Busca sequencial}
		\begin{itemize}
			\item \tcb{Elemento desejado:} 44.
		\end{itemize}
		\begin{center}
			\vetempty{0,1,2,3,4,5,6,7}\\
			\vet{22,\bf \bl 99,66,33,11,00,77,55}\\
		\end{center}
		\begin{itemize}
			\item \tcb{Comparação};
		\end{itemize}
	\end{frame}
	% ==========================================

	% ==========================================
	\begin{frame}[c]{Exemplo Busca sequencial}
		\begin{itemize}
			\item \tcb{Elemento desejado:} 44.
		\end{itemize}
		\begin{center}
			\vetempty{0,1,2,3,4,5,6,7}\\
			\vet{22,99,\bf \bl 66,33,11,00,77,55}\\
		\end{center}
		\begin{itemize}
			\item \tcb{Comparação};
		\end{itemize}
	\end{frame}
	% ==========================================

	% ==========================================
	\begin{frame}[c]{Exemplo Busca sequencial}
		\begin{itemize}
			\item \tcb{Elemento desejado:} 44.
		\end{itemize}
		\begin{center}
			\vetempty{0,1,2,3,4,5,6,7}\\
			\vet{22,99,66,\bf \bl 33,11,00,77,55}\\
		\end{center}
		\begin{itemize}
			\item \tcb{Comparação};
		\end{itemize}
	\end{frame}
	% ==========================================

	% ==========================================
	\begin{frame}[c]{Exemplo Busca sequencial}
		\begin{itemize}
			\item \tcb{Elemento desejado:} 44.
		\end{itemize}
		\begin{center}
			\vetempty{0,1,2,3,4,5,6,7}\\
			\vet{22,99,66,33,\bf \bl 11,00,77,55}\\
		\end{center}
		\begin{itemize}
			\item \tcb{Comparação};
		\end{itemize}
	\end{frame}
	% ==========================================

	% ==========================================
	\begin{frame}[c]{Exemplo Busca sequencial}
		\begin{itemize}
			\item \tcb{Elemento desejado:} 44.
		\end{itemize}
		\begin{center}
			\vetempty{0,1,2,3,4,5,6,7}\\
			\vet{22,99,66,33,11,\bf \bl 00,77,55}\\
		\end{center}
		\begin{itemize}
			\item \tcb{Comparação};
		\end{itemize}
	\end{frame}
	% ==========================================

	% ==========================================
	\begin{frame}[c]{Exemplo Busca sequencial}
		\begin{itemize}
			\item \tcb{Elemento desejado:} 44.
		\end{itemize}
		\begin{center}
			\vetempty{0,1,2,3,4,5,6,7}\\
			\vet{22,99,66,33,11,00,\bf \bl 77,55}\\
		\end{center}
		\begin{itemize}
			\item \tcb{Comparação};
		\end{itemize}
	\end{frame}
	% ==========================================

	% ==========================================
	\begin{frame}[c]{Exemplo Busca sequencial}
		\begin{itemize}
			\item \tcb{Elemento desejado:} 44.
		\end{itemize}
		\begin{center}
			\vetempty{0,1,2,3,4,5,6,7}\\
			\vet{22,99,66,33,11,00,77,\bf \bl 55}\\
		\end{center}
		\begin{itemize}
			\item \tcb{Comparação};
		\end{itemize}
	\end{frame}
	% ==========================================

	% ==========================================
	\begin{frame}[c]{Exemplo Busca sequencial}
		\begin{itemize}
			\item \tcb{Elemento desejado:} 44.
		\end{itemize}
		\begin{center}
			\vetempty{0,1,2,3,4,5,6,7}\\
			\vet{22,99,66,33,11,00,77,55}\\
		\end{center}
		\begin{itemize}
			\item \tcb{Comparação};
			\item \tcr{Não encontrou (posição: -1)}.
		\end{itemize}
	\end{frame}
	% ==========================================

	\section{Busca binária}

	\begin{frame}[c]{Busca binária}
		Em diversas aplicações reais, existe a necessidade de algoritmos
		eficientes. \tcr{Seria possível melhorar a eficiência do algoritmo
		de busca sequencial}?
	\end{frame}

	\begin{frame}[c]{Busca binária}
		Infelizmente, se os \tcr{elementos} estiverem \tcr{armazenados} em uma 
		\tcb{ordem aleatória} no vetor, \tcr{não} temos como melhorar o algoritmo
		de busca, pois precisamos verificar todos os elementos.
	\end{frame}
	
	\begin{frame}[c]{Busca binária}
		No entanto, se os elementos do vetor estiverem ordenados, podemos aplicar
		um algoritmo mais eficiente para realizar a busca. Trata-se do algoritmo
		\tcb{Busca Binária}.
	\end{frame}
	
	\begin{frame}[c]{Busca binária}
		Contexto:\\
		Jogo ``adivinhe o número'': pense em um número e diga o intervalo em que ele
		está. Para cada adivinhação de um oponente, apenas uma das três respostas
		podem ser dadas: \tcb{o número é maior, o número é menor ou acertou}. O vencedor
		é aquele que acertar o número com menos adivinhações.
	\end{frame}
	
	\begin{frame}[c]{Busca binária - conceito}
		\begin{enumerate}
			\item Testar o elemento que buscamos com o valor armazenado no \tcb{meio} do vetor;
			\item Se o elemento que buscamos for \tcb{menor} que o elemento do meio, sabemos que, se
			o elemento estiver presente no vetor, ele estará na primeira parte do vetor.
			\item Se for \tcb{maior}, estará na segunda parte do vetor;
			\item Se for \tcb{igual}, o elemento foi encontrado;
			\item Se concluirmos que o elemento está em uma das partes do vetor, 
				repetimos o procedimento descrito considerando apenas a parte restante.
		\end{enumerate}
	\end{frame}

	% ==========================================
	\begin{frame}[c]{Exemplo busca binária}
		\begin{itemize}
			\item \tcb{Elemento desejado:} 55.
		\end{itemize}
		\begin{center}
			\vetempty{0,1,2,3,4,5,6,7}\\
			\vet{00,11,22,33,55,66,77,99}\\
			\vetempty{\nd,\nd,\nd,\nd,\nd,\nd,\nd,\nd}\\
		\end{center}
		\begin{itemize}
			\item \tcb{Comparação};
			\item \tcr{Intervalo considerado (início e fim)};
		\end{itemize}
	\end{frame}
	% ==========================================
	
	% ==========================================
	\begin{frame}[c]{Exemplo busca binária}
		\begin{itemize}
			\item \tcb{Elemento desejado:} 55.
		\end{itemize}
		\begin{center}
			\vetempty{0,1,2,3,4,5,6,7}\\
			\vet{\re 00,11,22,\bl 33,55,66,77,\re 99}\\
			\vetempty{\huge \re i,\nd,\nd,\nd,\nd,\nd,\nd,\huge \re f}\\
		\end{center}
		\begin{itemize}
			\item \tcb{Comparação};
			\item \tcr{Intervalo considerado (início e fim)};
			\item \tcr{$\displaystyle \frac{(1+8)}{2} = 4$};
		\end{itemize}
	\end{frame}
	% ==========================================

	% ==========================================
	\begin{frame}[c]{Exemplo busca binária}
		\begin{itemize}
			\item \tcb{Elemento desejado:} 55.
		\end{itemize}
		\begin{center}
			\vetempty{0,1,2,3,4,5,6,7}\\
			\vet{00,11,22,33,\re 55,\bl 66,77,\re 99}\\
			\vetempty{\nd,\nd,\nd,\nd,\huge \re i,\nd,\nd,\huge \re f}\\
		\end{center}
		\begin{itemize}
			\item \tcb{Comparação};
			\item \tcr{Intervalo considerado (início e fim)};
			\item \tcr{$\displaystyle \frac{(5+8)}{2} = 6$};
		\end{itemize}
	\end{frame}
	% ==========================================
	
	% ==========================================
	\begin{frame}[c]{Exemplo busca binária}
		\begin{itemize}
			\item \tcb{Elemento desejado:} 55.
		\end{itemize}
		\begin{center}
			\vetempty{0,1,2,3,4,5,6,7}\\
			\vet{00,11,22,33,\re 55, 66,77, 99}\\
			\vetempty{\nd,\nd,\nd,\nd,\huge \re i f,\nd,\nd,\nd}\\
		\end{center}
		\begin{itemize}
			\item \tcb{Comparação};
			\item \tcr{Intervalo considerado (início e fim)};
			\item \tcr{$\displaystyle \frac{(5+5)}{2} = 5$};
		\end{itemize}
	\end{frame}
	% ==========================================

	% ==========================================
	\begin{frame}[c]{Exemplo busca binária}
		\begin{itemize}
			\item \tcb{Elemento desejado:} 55.
		\end{itemize}
		\begin{center}
			\vetempty{0,1,2,3,4,5,6,7}\\
			\vet{00,11,22,33,\re 55, 66,77, 99}\\
			\vetempty{\nd,\nd,\nd,\nd,\huge \re i f,\nd,\nd,\nd}\\
		\end{center}
		\begin{itemize}
			\item \tcb{Comparação};
			\item \tcr{Intervalo considerado (início e fim)};
			\item \tcr{Elemento encontrado na posição 5};
		\end{itemize}
	\end{frame}
	% ==========================================
	
	% ==========================================
	\begin{frame}[c]{Exemplo busca binária}
		\begin{itemize}
			\item \tcb{Elemento desejado:} 15.
		\end{itemize}
		\begin{center}
			\vetempty{0,1,2,3,4,5,6,7}\\
			\vet{00,11,22,33,55,66,77,99}\\
			\vetempty{\nd,\nd,\nd,\nd,\nd,\nd,\nd,\nd}\\
		\end{center}
		\begin{itemize}
			\item \tcb{Comparação};
			\item \tcr{Intervalo considerado (início e fim)};
		\end{itemize}
	\end{frame}
	% ==========================================
	
	% ==========================================
	\begin{frame}[c]{Exemplo busca binária}
		\begin{itemize}
			\item \tcb{Elemento desejado:} 15.
		\end{itemize}
		\begin{center}
			\vetempty{0,1,2,3,4,5,6,7}\\
			\vet{\re 00,11,22,\bl 33,55,66,77,\re 99}\\
			\vetempty{\huge \re i,\nd,\nd,\nd,\nd,\nd,\nd,\huge \re f}\\
		\end{center}
		\begin{itemize}
			\item \tcb{Comparação};
			\item \tcr{Intervalo considerado (início e fim)};
			\item \tcr{$\displaystyle \frac{(1+8)}{2} = 4$};
		\end{itemize}
	\end{frame}
	% ==========================================

	% ==========================================
	\begin{frame}[c]{Exemplo busca binária}
		\begin{itemize}
			\item \tcb{Elemento desejado:} 15.
		\end{itemize}
		\begin{center}
			\vetempty{0,1,2,3,4,5,6,7}\\
			\vet{\re 00,\bl 11,\re 22,33,55,66,77,\re 99}\\
			\vetempty{\huge \re i,\nd,\huge \re f,\nd,\nd,\nd,\nd,\nd}\\
		\end{center}
		\begin{itemize}
			\item \tcb{Comparação};
			\item \tcr{Intervalo considerado (início e fim)};
			\item \tcr{$\displaystyle \frac{(1+3)}{2} = 2$};
		\end{itemize}
	\end{frame}
	% ==========================================

	% ==========================================
	\begin{frame}[c]{Exemplo busca binária}
		\begin{itemize}
			\item \tcb{Elemento desejado:} 15.
		\end{itemize}
		\begin{center}
			\vetempty{0,1,2,3,4,5,6,7}\\
			\vet{00,11,\re 22,33,55,66,77,\re 99}\\
			\vetempty{\nd,\nd,\huge \re i \huge \re f,\nd,\nd,\nd,\nd,\nd}\\
		\end{center}
		\begin{itemize}
			\item \tcb{Comparação};
			\item \tcr{Intervalo considerado (início e fim)};
			\item \tcr{$\displaystyle \frac{(3+3)}{2} = 3$};
		\end{itemize}
	\end{frame}
	% ==========================================
	
	% ==========================================
	\begin{frame}[c]{Exemplo busca binária}
		\begin{itemize}
			\item \tcb{Elemento desejado:} 15.
		\end{itemize}
		\begin{center}
			\vetempty{0,1,2,3,4,5,6,7}\\
			\vet{00,11,22,33,55,66,77,\re 99}\\
			\vetempty{\nd,\nd,\huge \re i \huge \re f,\nd,\nd,\nd,\nd,\nd}\\
		\end{center}
		\begin{itemize}
			\item \tcb{Comparação};
			\item \tcr{Intervalo considerado (início e fim)};
			\item \tcr{Não encontrou (posição -1)};
		\end{itemize}
	\end{frame}
	% ==========================================

	\begin{frame}[fragile]{Busca Sequencial - Análise}

		\begin{minted}[
				frame=lines,
				framesep=2mm,
				baselinestretch=1.2,
				fontsize=\footnotesize,
				linenos
			]{c}
int busca(int * v, int n, int x){
	for(i=0; i<n; i++)
		if(v[i] == x)
			return(i)
	return(-1);
}
		\end{minted}
		\only<2-> \tcb{Quantas comparações} o algoritmo Busca Sequencial realiza?
		\begin{itemize}
			\item<3->{Melhor caso? \only<4->{$= 1$}}
			\item<3->{Pior caso? \only<5->{$\displaystyle = \sum_{i=0}^{n}{1} = n$}}
			\item<3->{Caso médio? \only<6->{$\displaystyle = \frac{n}{2}$}}
		\end{itemize}
	\end{frame}

	\begin{frame}[fragile]{Busca Binária}
		\begin{minted}[
				frame=lines,
				framesep=2mm,
				baselinestretch=1.2,
				fontsize=\footnotesize,
				linenos
			]{c}
int busca_binaria(int v[], int n, int elem){
  l = 0; r = n-1;

  while(l <= r){
    m = (l+r)/2;
    if(v[m] == elem) return(m);
    else if(v[m] < elem) l = m + 1;
    else if(v[m] > elem) r = m - 1;
  }
  return(-1);
}
		\end{minted}
	\end{frame}

	\begin{frame}[fragile]{Busca Binária Recursiva}
		\begin{minted}[
				frame=lines,
				framesep=2mm,
				baselinestretch=1.2,
				fontsize=\footnotesize,
				linenos
			]{c}
int busca_binaria_recur(int v[], int n, int elem, int l, int r){
  int m = (l+r)/2;

  if(l > r)
    return(-1);

  if(v[m] == elem)
    return(m);
  else if(v[m] < elem)
    return(busca_binaria_recur(v,n,elem,m+1,r));
  else if(v[m] > elem)
    return(busca_binaria_recur(v,n,elem,l,m-1));
}
		\end{minted}
	\end{frame}

	\begin{frame}{Busca Binária - Análise}
		Dado um inteiro $n$, quantas vezes é possível dividir $n$
		ao meio, até que o seu valor seja $1$ ou fique abaixo disso?

		\center $\displaystyle \frac{n}{2},\frac{n}{4},\frac{n}{8},\ldots,\frac{n}{2^k}$

	\end{frame}
	
	\begin{frame}{Busca Binária - Análise}
		\center $\displaystyle \frac{n}{2},\frac{n}{4},\frac{n}{8},\ldots,\frac{n}{2^k}$

		Para achar o valor de $k$ (número máximo de divisões):

		\center $\displaystyle \frac{n}{2^k} = 1$
		
		\center $\displaystyle k = \log_2{n}$
		
		\center $\displaystyle k = \ceil{\log_2{n}}$

		\center $\displaystyle \sum_{i=1}^{\ceil{\log_2{n}}}{1} = \log_2{n}$
	\end{frame}

	% ==========================================
	% \begin{frame}[fragile,c]{Implementação Recursiva Fatorial}
	%	\begin{minted}[frame=lines,framesep=2mm,baselinestretch=1.2,fontsize=\normalsize,linenos]{c}
	%	\end{minted}
	% \end{frame}
	% ==========================================

	% ==========================================
	% \begin{frame}
	% 	\frametitle{Título}
	% 	\begin{itemize}
	% 		\item Item bla bla bla
	% 	\end{itemize}
	%	  \begin{block}{Título de bloco}
	%		 	Conteúdo do bloco
	%   \end{block}
	% \end{frame}
	% ==========================================

	% ==========================================
	% Exemplo de frame com figura
	% \begin{frame}
	% 	\frametitle{Título}
	%	  \begin{figure}[h]
	%	  	\centering
	%	  	\includegraphics[width=10cm]{paht.png}
	%			\caption{Bla bla bla ...}
	%	  	\label{label}
	%	  \end{figure}
	% \end{frame}
	% ==========================================

	% ==========================================
	% Exemplo de frame com algoritmo
	% \begin{frame}[fragile]
	% 	\frametitle{Título}
	% 	\begin{lstlisting}[language=C]
	% 		#include <stdio.h>
	%
	% 		int main()
	% 		{
	% 		  int a[5];
	%
	% 		  a[0] = 30;
	% 		  a[1] = 25;
	% 		  a[2] = 5;
	% 		  a[3] = 100;
	% 		  a[4] = 1;
	%
	% 		  printf("%d\n",a[1]);
	% 		}
	% 	\end{lstlisting}
	% \end{frame}
	% ==========================================

	% ==========================================
	% Frame para referências
	%\begin{frame}[allowframebreaks]
	%	\frametitle{Referências utilizadas na apresentação}
	%	\bibliography{referencias}
	%\end{frame}
  % ==========================================

\end{document}
